\pagestyle{plain} % No headers, just page numbers

\chapter{\uppercase {Related Work}}

The proposed work draws on five areas of prior research: the theory of submodular
maximization, classical pruning methods, learning-based pruning methods, neural
combinatorial optimization and tree search, and automated feature engineering with large
language models. This chapter reviews existing work in each area and highlights the
limitations that motivate the contributions of this proposal.

\section{Submodular Maximization}

QuickPrune (Chapter~\ref{chap:quickprune}) builds on the theory of submodular
maximization, a framework that naturally models many graph CO problems such as maximum
coverage and influence maximization. \citet{nemhauser1978analysis} show that a greedy
algorithm achieves a $(1-1/e)$ approximation for monotone submodular maximization under a
cardinality constraint, a bound that is tight under standard complexity
assumptions~\citep{feige1998threshold}. \citet{vondrak2008optimal} extend this to
matroid constraints via the continuous greedy algorithm. QuickPrune generalizes beyond
standard submodularity by operating under $\gamma$-weak submodularity, a relaxation that
covers a broader class of objectives including influence maximization and sparse
regression~\citep{elenberg2018restricted}.

Several works address the scalability of submodular maximization through complementary
means. \citet{mirzasoleiman2015lazier} reduce the query complexity of greedy from $O(nk)$
to $O(n\log(1/\epsilon))$ via stochastic lazy evaluation, maintaining a
$(1-1/e-\epsilon)$ guarantee. \citet{badanidiyuru2014streaming} introduce a single-pass
streaming algorithm that retains a constant-factor approximation using $O(k/\epsilon)$
memory --- the streaming model that directly inspires QuickPrune's single-pass design. A
related line of work constructs \textit{coresets}: compact summaries of a dataset that
approximately preserve objective values and can be reused across multiple queries.
\citet{feldman2011unified} introduce a general coreset framework for a broad class of
shape-fitting and clustering problems, showing that small weighted subsets can substitute
for the full dataset with bounded error. \citet{mirzasoleiman2020coresets} extend this to
submodular maximization, constructing coresets that preserve the value of the optimal
solution up to a $(1-\epsilon)$ factor. QuickPrune's pruned universe serves a similar
role, but differs in two key respects: it operates under the value oracle model without
requiring reweighting of elements, and it supports a range of budgets simultaneously
rather than a single query type.

\section{Classical Approaches to Pruning}
Classical pruning methods for combinatorial optimization (CO) fall into two categories:
problem-specific heuristics and general-purpose algorithms. Problem-specific heuristics
exploit the structure of a single CO problem to discard unpromising elements, such as
degree-based candidate filtering for influence
maximization~\citep{chen2010scalable} and edge elimination rules for the traveling
salesman problem~\citep{sun2021generalization}. While effective within their respective
domains, these techniques do not transfer to other CO problems.

On the general-purpose side, \citet{zhou2017scaling} propose a pruning algorithm grounded
in submodularity to reduce the cost of submodular maximization under size constraints.
Their method constructs a pruned submodularity graph to identify and discard elements with
low marginal contributions. However, this approach processes elements sequentially, which
limits its scalability to large instances. More broadly, all existing classical pruning
methods share three limitations: they target a single instance of size-constrained
maximization, they provide no mechanism to leverage patterns across problem instances, and
they lack formal guarantees on the quality of the pruned ground set.

\section{Learning-Based Approaches to Pruning}
Learning-based methods take a data-driven approach to pruning, training models on solved
instances to identify elements that can be removed from the ground set.
These methods can be broadly organized into two categories: supervised approaches with
hand-crafted features and hybrid methods combining supervised and reinforcement learning.

\textbf{Supervised methods.} The earliest learning-based pruning methods train binary classifiers on small, pre-solved
instances to predict whether a vertex or edge belongs to an optimal solution. For Maximum
Clique Enumeration, \citet{lauri2019fine} apply logistic regression with hand-crafted
graph-theoretic features, and \citet{grassia2019learning, lauri2023learning} extend this
with multi-stage classifiers that iteratively reduce the search space. For the Traveling
Salesman Problem, \citet{fitzpatrick2021learning, sun2021generalization} train classifiers
to prune edges unlikely to belong to the optimal tour. \citet{tian2024combhelper} propose
COMBHelper, a GNN-based framework with knowledge distillation that still relies on
computationally expensive Fourier-based feature transformations and a problem-specific
boosting mechanism. A common limitation across all
these methods is their reliance on problem-specific feature engineering, which restricts
generalization to new CO problems.

\textbf{Hybrid supervised and reinforcement learning methods.}
Hybrid methods that combine supervised learning with reinforcement learning are more
problem-agnostic. GCOMB~\citep{manchanda2020gcomb} uses a probabilistic greedy
algorithm to generate training data, then applies a weighted-degree heuristic to prune
the ground set before training a Q-learning network on the reduced search space. However,
its initial pruning depends on a handcrafted degree heuristic that does not generalize
well across CO problems~\citep{ireland2022lense}. LeNSE~\citep{ireland2022lense} takes a
different approach: a GNN is first trained to generate embeddings for sampled subgraphs,
and a policy network then uses these embeddings to guide a local search toward subgraphs
likely to contain the optimal solution. A major drawback is its reliance on domain
expertise to categorize subgraphs, set class-specific sample counts, and tune
hyperparameters such as subgraph size and embedding dimensionality separately for each
dataset and problem.

Despite these advances, no existing learning-based pruning method simultaneously achieves
problem-agnostic design, automated feature engineering, and scalability to large
graph instances. Moreover, all existing methods are designed and evaluated exclusively
under size constraints, leaving other constraint types such as knapsack constraints
unaddressed.

\section{Neural Combinatorial Optimization and Tree Search}

A parallel line of work applies deep learning directly to CO problems in combination
with tree search. \citet{li2018combinatorial} combine a GCN with guided tree search in a
supervised setting, though their approach requires a large number of pre-solved instances
for training. \citet{abe2019solving} train a GCN using Monte Carlo Tree Search as a policy
improvement operator, following the AlphaGo Zero paradigm~\citep{silver2017mastering}.
\citet{xing2020graph} introduce a GNN-assisted MCTS method that constructs solutions to
the TSP by adding nodes successively. \citet{choo2022simulation} propose a fixed-width
tree search guided by a neural network policy. These methods focus on producing solutions
as close to optimal as possible, and are generally limited to graphs with hundreds to
thousands of nodes.

H-DeepPruner (Chapter~\ref{chap:hdeeppruner}) shares the GNN+MCTS design with this line
of work but is orthogonal in objective: rather than directly optimizing the CO objective,
it produces a reduced ground set for a separate downstream solver. This separation enables
scalability to large graphs and generalization across CO problems without retraining.

\section{Automated Feature Engineering via LLMs}
Automated feature engineering seeks to generate useful features from raw data without
human effort. Prior work with large language models (LLMs) has focused primarily on
tabular data, where new features are derived from existing columns. These approaches
either rely on manually designed feature-generation
rules~\citep{zhang2023openfe, zhang2024dynamic} or few-shot
prompts~\citep{han2024large}, and thus still depend on hand-crafted design choices.
Moreover, existing methods~\citep{nam2024optimized, hollmann2023large} suffer from error
propagation caused by simple feedback loops. Our approach is orthogonal to methods that directly generate algorithms using
LLMs~\citep{yang2023large, liu2024evolution, romera2024mathematical}, which o exhibit
degraded performance as graph size increases. Instead, LLM2Prune leverages
the reasoning capabilities of LLMs to generate features for a learned pruning algorithm,
keeping inference lightweight and transferable across instances without re-running the
search.

LLM-driven feature engineering is a natural fit for pruning, as LLMs can encode the kind
of domain knowledge that would otherwise require expert design of problem-specific
features. To the best of our knowledge, no prior work has applied LLM-driven feature
engineering to search space reduction for graph CO problems.